\documentclass{beamer}
\usepackage{tabularx}
\usepackage{graphicx}
\usepackage{multimedia}
\usepackage{hyperref}
\usepackage[norsk]{babel}
%Information to be included in the title page:
\title{Status automasjonssystemer på Flåte/Lukket merd}
\subtitle{Flåtestyring \- dødfisksystem \- Heimdall}
%\author{Anonymous}
\institute{Operational systems - ScaleAQ}
\date{\today}

\hypersetup{
  colorlinks, linkcolor=red
}

\begin{document}

\begin{frame}
  \frametitle{Barge Control}
    \begin{enumerate}
    \vspace{0pt} % <-- forces top alignment inside the minipage
    \item<1-> Lensepumpestyring
    \item<2-> Tanknivå 
    \item<3-> Hardwiret signaler til generator (start ved lekkasje), vannforing og F\&G-system
    \item<4-> Kommunikasjon til SiteControl over OPCUA
    \item<5-> Generatormonitorering (direktekommunikasjon i SiteControl)
    \item<6-> Kjøres på lab i Bergen (10.1.2.31) tilgjengelig over VPN (Global Protect)
    

    \end{enumerate}
    \includegraphics<1->[width=0.8\textwidth]{barge_dashboard.png}


\end{frame}


\begin{frame}
  \frametitle{Barge Control - utrulling til lokasjon}
  \begin{enumerate}
    \vspace{0pt} % <-- forces top alignment inside the minipage
    \item<1-> \href{https://steinsvik.sharepoint.com/sites/ScaleAQService}{ScaleAQ Service $\,\to\,$ Barge Control 2.0}
    \item<2-> Tekniker bruker software WAGOupload. Enkeltfil som inneholder codesys-applikasjon med config og fw
    \item<2-> Enkeltfilen (.appload) som brukes av WAGOupload må genereres fra kompilert codesys-applikasjon.
    \item<4-> Prosedyre står beskrevet i `Commissioning guide' i ScaleAQ Service.
    \end{enumerate}

    \includegraphics<2>[width=0.8\textwidth]{WAGOupload.png}
\end{frame}

\begin{frame}
  \frametitle{Barge Control - vedlikehold}
  \begin{enumerate}
    \vspace{0pt} % <-- forces top alignment inside the minipage
    \item<1-> Utrullet på 10+ lokasjoner
    \item<2-> Skal i utgangspunktet ikke være mye oppdateringer
    \item<3-> Teknikere er kjent med utrullingsprosedyre og har vært lite behov for hjelp til dette
    \item<4-> Kan være behov for å ha en mer automatisert måte og oppdatere program til PLS'er på. F.eks ´Wago Device Sphere' eller lignende. Begynner å bli tungvint med manuelle oppdateringer via feedstation VM
    \end{enumerate}

\end{frame}

\begin{frame}
  \frametitle{Automatisert dødfisksystem - prosessflyt}
\begin{minipage}[t]{0.5\textwidth}
\scriptsize 
    \begin{enumerate}
    \vspace{0pt} % <-- forces top alignment inside the minipage
    \item<1-> Død fisk samles i bunn av merd
    \item<2-> Fisk blir ledet til flåte ved hjelp av trykkluft
    \item<3-> Stunnerenhet(myndighetspålagt ved remotestyring av prosess)
    \item<4-> Beveger seg gjennom LiftUp-enhet på et VSD-styrt rullebånd
    \item<5-> Vaskepumpe
    \item<6-> Kameraenhet / visionsystem
    \item <7-> Avstandssensor (ultrasonisk)
    \item <8-> Blir sluppet ned i grovkvern
    \item <9-> Blir lagret i tanker og sirkulert av ensilasjesystemet
    \item <10-> Automatisk håndtering av syredosering
    \end{enumerate}
\end{minipage}%
\hfill%
\begin{minipage}[t]{0.5\textwidth}
    \vspace{0pt} % top alignment for images
    \includegraphics<1>[width=0.8\textwidth]{deadfishpen.png}
    \includegraphics<2>[width=0.7\textwidth]{airliftpumppres.png}
    \includegraphics<3>[width=0.8\textwidth]{liftup2.jpg}
    \includegraphics<4-6>[width=0.8\textwidth]{liftupmodel.png}
    \includegraphics<7>[width=0.8\textwidth]{fishpresent.png}
    %
    \includegraphics<7>[width=0.8\textwidth]{sick.png}
    \includegraphics<8>[width=0.8\textwidth]{pregrinder.png}
    \includegraphics<9-10>[width=0.8\textwidth]{ensilage.png}

\end{minipage}
\end{frame}

\begin{frame}
  \frametitle{Dødfisksystem - nettverkstopologi}
    \includegraphics<1->[width=0.8\textwidth]{networktopo.drawio.png}


\end{frame}

\begin{frame}
  \frametitle{Dødfisksystem - manuell modus}
  \begin{enumerate}
    \vspace{0pt} % <-- forces top alignment inside the minipage
    \item<1-> Kan starte delprosesser manuelt.
    \item<2-> Ikke avhengig av at alle underprosesser er tilgjengelig for styresystemte (kompressor, ensilasje, stunner etc.)
    \item<3-> Typisk use case er:
    \begin{enumerate}
      \item<4-> Se hvilke undersystemer som er \texttt{available} (tilgjengelig og uten feilmeldinger)
      \item <5-> Start kompressor - se at statussignal endres til \texttt{running}
      \item <6-> Start rullebånd - verifiser at vi får \texttt{running} signal fra VSD
      \item <7-> Send signal til liftup om hvilken merd vi skal pumpe fra (1,2,3,...) og verifiser at vi får \texttt{pen x active} signal fra system.
      \item <8-> send \texttt{dead fish active} signal til ensilasjesystem.
      \item <9-> Stopp alle delprosesser når man er ferdig
    \end{enumerate}
    \end{enumerate}

\end{frame}

\begin{frame}
  \frametitle{Dødfisksystem - semiauto modus}
  \begin{enumerate}
    \vspace{0pt} % <-- forces top alignment inside the minipage
    \item<1-> Nødvendige delsystemer må være \texttt{available} og \texttt{enabled}. Dette vises i \texttt{status} strukturen for hele liftup-objektet (\texttt{DeadFishUnit}).
    \item<2-> Setter systemet i \texttt{semi-auto} modus.
    \item<3-> Aktiver \texttt{penStartPumping} (UINT)
    \item<4-> Verifiser at nødvendige delsystemer starter opp. Se at vi får tilsvarende \texttt{penRunning}-signal tilbake fra status-strukturen.
    \item<5-> Gi ny verdi til \texttt{penStartPumping} for pumping på en annen merd.
    \item<6-> Stopp systemet ved å sende \texttt{penStopPumping}.
    \end{enumerate}
  \end{frame}



    \begin{frame}
  \frametitle{Dødfisksystem - auto modus}
  \begin{enumerate}
    \vspace{0pt} % <-- forces top alignment inside the minipage
    \item<1-> Skriv config-verdier for hvilke merder som er aktive, hvor lenge hver enkelt merd skal pumpes fra og når man starter sekvensen på dagen og hvor hyppig.
    \item<2-> Programmet tar hånd om at alle delsystemer er tilgjengelige, og at ting starter opp i riktig sekvens.
    \item<3-> Denne modusen er ikke ferdig programmert - eller testet.
    \item<4-> Tanken er at i første omgang kan dette kun startes i manuel / semi-auto modus. Når systemet er gjennomtestet nok og vi føler oss trygg på at dette kan kjøres schedule-basert så implementerer vi denne funksjonen.
    \end{enumerate}
\end{frame}

    \begin{frame}
  \frametitle{Dødfisksystem - status}
  \begin{enumerate}
    \vspace{0pt} % <-- forces top alignment inside the minipage
    \item<1-> Program er utrullet på enilasje-PLS på Buholmen.
    \item<2-> Kommunikasjon er opprettet mot LiftUp-system.
    \item<3-> Styremodul for kompressor er ikke installert eller testet.
    \item<4-> Vi har testet remote manuell start av systemene. (semi-auto kunne ikke testes pga. manglende oppsett i LiftUp-enhet). Programmet var avhengig av at stunner og flusher var oppe å gå, noe det ikke var under test-tidspunkt.
    \item<5-> Testrapport legges i Confluence-siden sammen med funksjonsbeskrivelser.
    \item<6-> Når SiteControl skal begynne og integrere mot PLS må noen ta hånd om at PLS-program kjøres i simuleringsmodus, og bistå teamet med eventuell inormasjon og oppdateringer.
    \item<7-> PLS kjører på 10.1.2.35 på labben i bergen. Liftup-simulering (modbus-server) kjøres i dockerfil på 10.1.2.101
    \end{enumerate}
\end{frame}



\begin{frame}
  \frametitle{CAS - Heimdall}
\begin{minipage}[t]{0.5\textwidth}
\scriptsize 
    \begin{enumerate}
    \vspace{0pt} % <-- forces top alignment inside the minipage
    \item<1-> Lukket system i sjø
    \item<2-> Tre VSD-styrte neddykkede inntaktspumper med tilhørende PLS fra leverandør.
    \item<3-> Hydraulikkstyrt regulerbar utløpsventil. Estimerer posisjon ved å starte/stoppe hydraulikkpumpe, ventiler og måling av oljeflow.
    \item<4-> Tre stk. oksygeneringssystem som henger på tilhørende inntaksdelstrøm
    \item<6-> Diverse IO for sensorikk og ventiler/aktuatorer
    \item<7-> Modbuskommunikasjon mot 3 stk RIO plassert på merdkant
    \item<8-> OPC-kommunikasjon mot 3 stk. OPC-servere for lesing/styring fra tilhørende VSD (pumpeinnstrømning)
    \end{enumerate}
\end{minipage}%
\hfill%
\begin{minipage}[t]{0.5\textwidth}
    \vspace{0pt} % top alignment for images
    \includegraphics<1-2>[width=1\textwidth]{cas.png}
    \includegraphics<3,6,8>[width=1\textwidth]{cas_subsea.png}
    \includegraphics<4>[width=1\textwidth]{cas_top_oxy.png}
    \includegraphics<7>[width=1\textwidth]{magnus.png}
    \includegraphics<7>[width=1\textwidth,trim=0 0 0 -3cm]{scale_lin.png}
\end{minipage}
\end{frame}














\end{document}